\chapter{导数：初步}\label{ch:g11:deriv}

函数变化得有多快？某一点处的答案是
\emph{导数}，应用数学中最有用的一个数。
本章由平均变化率出发构造导数，
计算基本函数的导数，并利用其符号读出
函数的变化情况。本章内容将在
\cref{ch:g12:deriv}中得到推广和深化。

\section{平均变化率}

\begin{definition}[平均变化率]\label{def:g11:deriv:rate}
设 $f$ 在区间 $I$ 上有定义，$a \neq b$ 为 $I$ 中的两个数。
\emph{平均变化率}\index{平均变化率}是 $f$ 在
$a$ 与 $b$ 之间的变化量之比，即
\[
\frac{f(b) - f(a)}{b - a}.
\]
它是曲线上过两点的\emph{割线}\index{割线}的斜率。
这两点为 $\bigl(a, f(a)\bigr)$ 与 $\bigl(b, f(b)\bigr)$。
\end{definition}

\begin{example}
一辆汽车在下午 2 点至 4 点间行驶了 $140$ 千米，
其平均速度为 $\frac{140}{2} = 70$ 千米/时。然而速度表在每一时刻都显示一个
\emph{瞬时}速度 --- 这正是
导数所要精确刻画的概念。
\end{example}

\section{一点处的导数}

\begin{definition}[一点处的导数]\label{def:g11:deriv:derivative}
设 $f$ 在包含 $a$ 的区间上有定义。对 $h \neq 0$，
作出 $a$ 与 $a + h$ 之间的平均变化率：
\[
r(h) = \frac{f(a+h) - f(a)}{h}.
\]
若 $r(h)$ 随 $h$ 任意接近
$0$ 时趋于一个确定的数，则称 $f$ \emph{在 $a$ 处可导}\index{导数}\index{可导}；
此数即为\emph{$f$ 在 $a$ 处的导数}，记作 $f'(a)$：
\[
f'(a) = \lim_{h \to 0} \frac{f(a+h) - f(a)}{h}.
\]
\end{definition}

\begin{example}\label{ex:g11:deriv:atpoint}
设 $f(x) = x^2$ 且 $a = 1$。对 $h \neq 0$：
\[
\frac{f(1+h) - f(1)}{h}
= \frac{(1+h)^2 - 1}{h}
= \frac{1 + 2h + h^2 - 1}{h}
= \frac{h(2 + h)}{h}
= 2 + h .
\]
当 $h$ 趋于 $0$ 时，上式趋于 $2$，故函数 $x^2$
在 $1$ 处可导，且 $f'(1) = 2$。注意策略：\emph{化简商式，
直到 $h$ 不再出现在分母中}，再令
$h \to 0$。
\end{example}

\begin{example}[不可导的点]
设 $f(x) = \abs{x}$ 且 $a = 0$。则
$\frac{f(h) - f(0)}{h} = \frac{\abs h}{h}$，其值为 $1$（当 $h > 0$ 时），
而 $-1$（当 $h < 0$ 时），并无单一数被趋近，故 $\abs x$ 在
$0$ 处不可导。其曲线在该点有一个\emph{角点}。
\end{example}

\section{切线}

\begin{definition}[切线]\label{def:g11:deriv:tangent}
若 $f$ 在 $a$ 处可导，则\emph{切线}\index{切线}
为曲线在点 $A = \bigl(a, f(a)\bigr)$ 处的直线，它经过 $A$ 且
斜率为 $f'(a)$，其方程为：
\[
y = f(a) + f'(a)(x - a).
\]
\end{definition}

切线是过 $A$ 的割线的极限位置：当
$h \to 0$ 时，第二点 $\bigl(a+h, f(a+h)\bigr)$ 沿
曲线向 $A$ 滑动，割线的斜率 $\frac{f(a+h)-f(a)}{h}$
趋近于 $f'(a)$。

\begin{omfigure}
\begin{tikzpicture}
\begin{axis}[omaxis,
  width=8.8cm, height=6cm,
  xmin=-0.4, xmax=3.2, ymin=-0.6, ymax=4.3,
  xtick={1,2.2}, xticklabels={$a$,$a+h$}, ytick=\empty]
  \addplot[omProp, thick, domain=-0.1:2.9] {1 + 1.6*(x-1)};
  \addplot[omThm, thick, domain=-0.35:2.6] {1 + 1*(x-1)};
  \addplot[omDef, thick, domain=-0.4:2.75] {x^2/2 + 0.5};
  \fill (axis cs:1,1) circle (1.5pt) node[above left, font=\small] {$A$};
  \fill (axis cs:2.2,2.92) circle (1.5pt)
    node[above left, font=\small] {$B$};
  \node[omProp, font=\small, anchor=west] at (axis cs:2.62,3.75)
    {割线};
  \node[omThm, font=\small, anchor=west] at (axis cs:2.62,2.45)
    {切线};
\end{axis}
\end{tikzpicture}

{\small 过 $A$ 与 $B = \bigl(a+h, f(a+h)\bigr)$ 的割线，其斜率为
$\frac{f(a+h)-f(a)}{h}$。当 $B$ 沿曲线滑向 $A$ 时，割线
转为切线，其斜率为 $f'(a)$。}
\end{omfigure}

\begin{example}\label{ex:g11:deriv:tangenteq}
对 $f(x) = x^2$，取 $a = 1$：$f(1) = 1$ 且 $f'(1) = 2$
（见\cref{ex:g11:deriv:atpoint}），故切线为
\[
y = 1 + 2(x - 1) = 2x - 1 .
\]
这正是在\cref{exo:g11:quad:11}中用纯代数方法求得的直线 $d_2$。
\end{example}

\section{基本函数的导数}

\begin{proposition}\label{prop:g11:deriv:refderivatives}
在各自的可导定义域上：
\[
(c)' = 0, \quad
(x)' = 1, \quad
(x^2)' = 2x, \quad
(x^3)' = 3x^2, \quad
\left(\frac1x\right)' = -\frac{1}{x^2}, \quad
(\sqrt x\,)' = \frac{1}{2\sqrt x} \ (x > 0).
\]
更一般地，$(x^n)' = n\,x^{n-1}$ 对每个整数 $n \geq 1$ 成立。
\end{proposition}

\begin{proof}
各证明均沿用\cref{ex:g11:deriv:atpoint}中的思路：先化简
变化率，再令 $h \to 0$。取定定义域中的 $a$。

\emph{常值函数与恒等函数。}$\frac{c - c}{h} = 0$ 与
$\frac{(a+h) - a}{h} = 1$ 对所有 $h$ 成立，极限分别为 $0$ 与 $1$。

\emph{平方。}$\frac{(a+h)^2 - a^2}{h} = \frac{2ah + h^2}{h} = 2a + h
\to 2a$。

\emph{立方。}$(a+h)^3 = a^3 + 3a^2h + 3ah^2 + h^3$，于是变化率为
$3a^2 + 3ah + h^2 \to 3a^2$。

\emph{倒数。}对 $a \neq 0$ 及充分小的 $h$，使得
$a + h \neq 0$：
\[
\frac{1}{h}\left(\frac{1}{a+h} - \frac{1}{a}\right)
= \frac{1}{h} \cdot \frac{a - (a+h)}{a(a+h)}
= \frac{-1}{a(a+h)}
\xrightarrow[h \to 0]{} -\frac{1}{a^2}.
\]

\emph{平方根。}对 $a > 0$，乘以共轭式：
\[
\frac{\sqrt{a+h} - \sqrt a}{h}
= \frac{(a + h) - a}{h\,(\sqrt{a+h} + \sqrt a)}
= \frac{1}{\sqrt{a+h} + \sqrt a}
\xrightarrow[h \to 0]{} \frac{1}{2\sqrt a}.
\]

\emph{一般幂。}$2a$、$3a^2$ 所示的规律仍然成立：展开
$(a+h)^n$ 中，$a^n + n\,a^{n-1}h$ 之后的各项都含 $h^2$，故变化率
趋于 $n\,a^{n-1}$。（完整的归纳法证明在
\cref{ch:g12:deriv}中给出。）
\end{proof}

\section{求导法则}

\begin{theorem}[运算法则]\label{thm:g11:deriv:rules}
设 $u$ 与 $v$ 在 $I$ 上可导，$\lambda \in \R$。则
$u + v$、$\lambda u$、$uv$ 在 $I$ 上可导，
$\frac{u}{v}$ 在使 $v \neq 0$ 的点处亦可导，且
\[
(u+v)' = u' + v', \qquad
(\lambda u)' = \lambda u', \qquad
(uv)' = u'v + uv', \qquad
\left(\frac{u}{v}\right)' = \frac{u'v - uv'}{v^2}.
\]
\end{theorem}

\begin{proof}[前三条法则的证明]
取定 $a \in I$，并记 $\Delta u = u(a+h) - u(a)$，
$\Delta v = v(a+h) - v(a)$。

\emph{和。}$u + v$ 的变化率为
$\frac{\Delta u + \Delta v}{h} = \frac{\Delta u}{h} + \frac{\Delta v}{h}
\to u'(a) + v'(a)$。

\emph{数乘。}$\lambda u$ 的变化率为
$\lambda \frac{\Delta u}{h} \to \lambda u'(a)$。

\emph{乘积。}加上并减去混合项 $u(a+h)\,v(a)$：
\begin{align*}
\frac{u(a+h)v(a+h) - u(a)v(a)}{h}
&= \frac{u(a+h)v(a+h) - u(a+h)v(a) + u(a+h)v(a) - u(a)v(a)}{h}\\
&= u(a+h)\,\frac{\Delta v}{h} + v(a)\,\frac{\Delta u}{h}.
\end{align*}
当 $h \to 0$ 时，$\frac{\Delta v}{h} \to v'(a)$，$\frac{\Delta u}{h} \to
u'(a)$，且 $u(a+h) \to u(a)$（其变化率有极限，故
$\Delta u = h \cdot \frac{\Delta u}{h} \to 0$）。于是变化率趋于
$u(a)v'(a) + v(a)u'(a)$。

商的求导法则同理可证；它\emph{在本阶段
不加证明地承认}，并在\cref{ch:g12:deriv}中证明。
\end{proof}

\begin{example}\label{ex:g11:deriv:computations}
\emph{多项式。}$f(x) = 3x^3 - 5x^2 + 7x - 2$：
\[
f'(x) = 3 \times 3x^2 - 5 \times 2x + 7 = 9x^2 - 10x + 7 .
\]

\emph{商。}$g(x) = \frac{2x+1}{x^2+1}$：其中 $u = 2x + 1$，$v = x^2
+ 1$，
\[
g'(x) = \frac{2(x^2+1) - (2x+1)\,2x}{(x^2+1)^2}
= \frac{-2x^2 - 2x + 2}{(x^2+1)^2}.
\]

\emph{乘积。}$k(x) = x\sqrt x$，当 $x > 0$ 时：
$k'(x) = 1 \times \sqrt x + x \times \frac{1}{2\sqrt x}
= \sqrt x + \frac{\sqrt x}{2} = \frac{3\sqrt x}{2}$。
\end{example}

\section{导数的符号与单调性}

\begin{theorem}[导数与单调性]\label{thm:g11:deriv:variations}
设 $f$ 在区间 $I$ 上可导。
\begin{enumerate}
  \item 若 $f' \geq 0$ 在 $I$ 上，则 $f$ 在 $I$ 上递增。
  \item 若 $f' \leq 0$ 在 $I$ 上，则 $f$ 在 $I$ 上递减。
  \item 若 $f' > 0$ 在 $I$ 上（除有限个使其为零的
        点外），则 $f$ 在 $I$ 上严格递增。
\end{enumerate}
\end{theorem}

\begin{proof}
直观上是明显的 --- 切线全部朝上的曲线必然上升 ---
但严格的证明需要用到中值定理。该结论
\emph{在本阶段不加证明地承认}；见\cref{thm:g12:deriv:variations}。
\end{proof}

\begin{proposition}[极值]\label{prop:g11:deriv:extrema}
若 $f'$ 在 $a$ 处由 $+$ 变为 $-$，则 $f$ 在 $a$ 处取得
极大值；若由 $-$ 变为 $+$，则为极小值。此点处 $f'(a) = 0$，
切线水平。
\end{proposition}

\begin{proof}
若 $f' \geq 0$ 在以 $a$ 为右端点的区间上，且 $f' \leq 0$
在以 $a$ 为左端点的区间上，则 $f$ 递增至 $f(a)$，之后递减。
由\cref{thm:g11:deriv:variations}可知：$f(a)$ 是附近的最大值。由于 $f'$
在 $a$ 处变号，且在该点有定义，故 $f'(a) = 0$。
\end{proof}

\begin{method}[研究函数]\label{met:g11:deriv:study}
\begin{enumerate}
  \item 求 $f'(x)$ 并化简 --- 最好化为因式分解形式。
  \item 确定 $f'(x)$ 在各区间上的符号（若
        $f'$ 为二次多项式，可用\cref{thm:g11:quad:sign}）。
  \item 将结果记录于\emph{变化表}中：第一行写出
        $f'$ 的符号，第二行用箭头表示 $f$ 的变化，并标出
        $f$ 在各转折点处的值。
  \item 得出结论：极值、方程 $f(x) = k$ 的解的个数、不等式。
\end{enumerate}
\end{method}

\begin{example}\label{ex:g11:deriv:study}
研究函数 $f(x) = x^3 - 3x + 1$ 在 $\R$ 上的性质。首先，
$f'(x) = 3x^2 - 3 = 3(x-1)(x+1)$，为二次多项式，根为 $-1$ 与 $1$，
在两根之外为正。于是 $f$ 在 $\intoc{-\infty}{-1}$ 上递增，
在 $\intcc{-1}{1}$ 上递减，在 $\intco{1}{+\infty}$ 上递增，并有
极大值 $f(-1) = -1 + 3 + 1 = 3$ 与极小值
$f(1) = 1 - 3 + 1 = -1$。
\end{example}

\begin{omfigure}
\begin{tikzpicture}
\begin{axis}[omaxis,
  width=8.4cm, height=5.8cm,
  xmin=-2.6, xmax=2.6, ymin=-3.6, ymax=4.4,
  xtick={-1,1}, ytick={-1,3}]
  \addplot[omProp, thick, domain=-1.7:-0.3] {3};
  \addplot[omProp, thick, domain=0.3:1.7] {-1};
  \addplot[omDef, thick, domain=-2.25:2.25, samples=90] {x^3 - 3*x + 1};
  \fill (axis cs:-1,3) circle (1.5pt);
  \fill (axis cs:1,-1) circle (1.5pt);
\end{axis}
\end{tikzpicture}

{\small 函数 $f(x) = x^3 - 3x + 1$ 的图像
（见\cref{ex:g11:deriv:study}）：水平切线（橙色）位于
极大值点 $(-1, 3)$ 与极小值点 $(1, -1)$ 处，
$f'$ 在这些点处变号。}
\end{omfigure}

\begin{example}[优化问题]\label{ex:g11:deriv:box}
一个无盖的盒子由一张 $12 \times 12$ cm 的方纸板制成：在每个角上
剪去边长为 $x$ 的正方形，再折起。其体积为
$V(x) = x(12 - 2x)^2$，$0 < x < 6$。展开得
$V(x) = 4x^3 - 48x^2 + 144x$，于是
\[
V'(x) = 12x^2 - 96x + 144 = 12(x^2 - 8x + 12) = 12(x - 2)(x - 6).
\]
在 $\intoo{0}{6}$ 上，因子 $x - 6$ 为负，故 $V'(x)$ 与
$-(x-2)$ 同号：在 $2$ 之前为正，之后为负。体积在
$x = 2$ cm 时最大，此时 $V(2) = 2 \times 8^2 = 128$ cm$^3$。
\end{example}

\section{练习}

\begin{exercise}[$\star$]\label{exo:g11:deriv:1}
利用定义（先求变化率，再令 $h \to 0$），
求 $f'(2)$（$f(x) = x^2$），再求 $g'(1)$（$g(x) = x^2 + 3x$）。
\end{exercise}

\begin{exercise}[$\star$]\label{exo:g11:deriv:2}
求下列函数的导数：
\[
f(x) = 4x^3 - 2x^2 + x - 7, \qquad
g(x) = 5\sqrt{x} + \frac{3}{x}, \qquad
h(x) = (2x+1)(x^2 - 3) .
\]
\end{exercise}

\begin{exercise}[$\star$]\label{exo:g11:deriv:3}
求函数 $f$ 的图像在横坐标为
$a$ 的点处的切线方程：
\[
f(x) = x^2 - 3x, \ a = 2; \qquad
f(x) = \frac1x, \ a = 1; \qquad
f(x) = \sqrt x, \ a = 4 .
\]
\end{exercise}

\begin{exercise}[$\star$]\label{exo:g11:deriv:4}
求下列分式的导数：
\[
f(x) = \frac{x+2}{x-1}, \qquad
g(x) = \frac{x^2}{x^2+1}, \qquad
h(x) = \frac{1}{x^2 + x + 1} .
\]
\end{exercise}

\begin{exercise}[$\star\star$]\label{exo:g11:deriv:5}
讨论 $f(x) = x^3 - 6x^2 + 9x - 2$ 在 $\R$ 上的变化情况（画出变化表，
并求局部极值）。
\end{exercise}

\begin{exercise}[$\star\star$]\label{exo:g11:deriv:6}
讨论 $f(x) = \frac{x^2 + 3}{x + 1}$ 在
$\intoo{-1}{+\infty}$ 上的变化情况。
\end{exercise}

\begin{exercise}[$\star\star$]\label{exo:g11:deriv:7}
设 $f(x) = x^2 - 4x + 5$。
\begin{enumerate}
  \item 求曲线上切线为水平直线的点。
  \item 求曲线上切线与直线
        $y = 2x + 1$ 平行的点。
\end{enumerate}
\end{exercise}

\begin{exercise}[$\star\star$]\label{exo:g11:deriv:8}
用 $60$ m 长的篱笆靠墙围成一个矩形场地
（篱笆围三边）。将面积表示为宽 $x$ 的函数，并利用
导数求面积最大时的尺寸。（可与
\cref{ch:g11:quad}中的顶点法比较。）
\end{exercise}

\begin{exercise}[$\star\star$]\label{exo:g11:deriv:9}
证明：对所有 $x > 0$，有 $\sqrt{x} \leq \frac{x + 1}{2}$，
通过研究函数 $f(x) = \frac{x+1}{2} - \sqrt x$ 在 $\intoo{0}{+\infty}$ 上的情况来证明。
\end{exercise}

\begin{exercise}[$\star\star$]\label{exo:g11:deriv:10}
一个圆柱形罐头须能容纳 $500$ cm$^3$。其表面积（两个圆盘与
侧面之和）为 $S = 2\pi r^2 + 2\pi r h$，且 $\pi r^2 h = 500$。将 $S$
表示为仅关于 $r$ 的函数，并证明 $S'(r) = 4\pi r -
\frac{1000}{r^2}$。由此求使表面积最小的半径，
并用精确形式表示。
\end{exercise}

\begin{exercise}[$\star\star\star$]\label{exo:g11:deriv:11}
有几条切线过抛物线 $y = x^2$ 外的一点
$P(1, -3)$？求出它们的方程。（设 $a$ 为切点的
横坐标，并由“$a$ 处的切线过 $P$”列出条件。）
\end{exercise}

\begin{exercise}[$\star\star\star$]\label{exo:g11:deriv:12}
设 $f(x) = x^3 - 3x + 1$（见\cref{ex:g11:deriv:study}）。利用
变化表，求方程
$f(x) = k$ 的解的个数，按实数 $k$ 的取值而定。
\end{exercise}

\section{问题：切线的三重作用}

\begin{problem}[{周末问题 --- 抛物面镜得证，
牛顿求根机与圆木中最坚固的梁}]
\label{pb:g11:deriv:1}
导数画切线，听起来只是一项谦逊的本领。
本题展示切线同时承担三项任务：其一，
证明了
\cref{pb:g11:quad:1}中留下的车灯定理（与抛物线对称轴平行的每条光线
经焦点反射）；其二，支撑当今常用的最快
解方程算法 --- 在每台计算器与神经网络之中 ---
其三是
求出圆木所能截出的最坚固矩形梁。

\medskip
\noindent\textbf{第一部分 --- 切线运算训练。}
\begin{enumerate}
  \item 求 $f(x) = 3x^2 - 5x + 1$ 与
        $g(x) = x^3 - 12x$ 的导数
        （由\cref{thm:g11:deriv:rules}）；并由定义求 $\left(x^2
        \right)'$ 在 $x = 1$ 处的导数（先求平均变化率，
        再令 $h \to 0$）。
  \item 求曲线 $y = x^3 - 12x$ 在 $a = 2$ 处的切线。该切线有何
        特殊之处，这说明了什么
        （见\cref{prop:g11:deriv:extrema}）？
  \item 求抛物线 $y = x^2$
        在点 $(a, a^2)$ 处的一般切线方程：
        \[
        y = 2ax - a^2 .
        \]
  \item 该切线与 $y$ 轴交于何处？由此得出
        制图员的作图法则：要作抛物线上的点 $P$
        （高为 $h$）处的切线，将 $P$ 与
        轴上原点下方深度为\,\dots\ 的点相连。请叙述此法则。
  \item 作出
        $g(x) = x^3 - 12x$ 的完整变化表（$g'$ 的符号、局部极值及
        其数值，见\cref{met:g11:deriv:study}）。
\end{enumerate}

\medskip
\noindent\textbf{第二部分 --- 镜面定理。}
设 $P(a, a^2)$ 为抛物线 $y = x^2$ 上的点（取
$a \neq 0$），$F\left(0, \frac14\right)$ 为焦点，$T$ 为
过 $P$ 的切线与 $y$ 轴的交点（第
4 小题）。由\cref{pb:g11:quad:1}可知
$PF = a^2 + \frac14$。
\begin{enumerate}[resume]
  \item 求距离 $FT$。
  \item 由此得出三角形 $FPT$ 在 $F$ 处为等腰三角形。
        因此哪两个角相等？
  \item 过 $P$ 且平行于对称轴的入射光线，
        与直线 $(TF)$ 平行（两者均为竖直方向）。利用
        内错角，证明过 $P$ 的切线与入射竖直光线
        所成的\emph{角}，与线段 $[PF]$ 所成的角
        相等。
  \item 物理：光线在曲线上的反射与其在切线上的反射相同，
        出射角等于入射角（即
        反射定律）。由此得出\emph{镜面定理}：每条
        与对称轴平行的光线经反射后都过焦点 ---
        反过来，焦点处的灯泡发出
        平行光。抛物面镜与车灯
        （见\cref{pb:g11:quad:1}）现在都已成为定理。
  \item 取 $a = 1$，对等腰三角形进行数值验证：
        计算 $T$、$FP$ 与 $FT$。
\end{enumerate}

\medskip
\noindent\textbf{第三部分 --- 牛顿的“机器”。}
为解方程 $f(x) = 0$，牛顿提出：先取一个猜测值
$x_0$，沿 $x_0$ 处的切线向下作到坐标轴，
其 $x$ 截距即为下一个猜测值。
\begin{enumerate}[resume]
  \item 推导该方法的公式：
        \[
        x_{n + 1} = x_n - \frac{f(x_n)}{f'(x_n)} .
        \]
  \item 对 $f(x) = x^2 - 2$ 从 $x_0 = 1$ 出发运行该方法：计算
        $x_1$、$x_2$、$x_3$，用精确分数表示。哪一个
        两千年前的数列 --- 以及
        \cref{pb:g10:functions:1}中的哪台机器 --- 你刚刚
        重建了？请证明这一巧合背后的恒等式：
        \[
        x - \frac{x^2 - 2}{2x} = \frac12\left(x +
        \frac2x\right).
        \]
  \item 解方程 $x^3 = 100$（即
        \cref{pb:g10:reffunc:1}中的储蓄交叉问题）：对
        $f(x) = x^3 - 100$ 以 $x_0 = 5$ 为初值，做两步牛顿迭代（保留四位小数），并
        与计算器的 $\sqrt[3]{100}$ 作比较。
  \item 破坏实验：尝试对
        $f(x) = x^3 - 12x$ 在 $x_0 = 2$ 处启动该方法，会出什么问题？
        为什么会这样（见第 2 题）？并指出实际使用中的注意事项。
  \item 在第 12 题中，正确的小数位数大致按
        $1 \to 3 \to 6$ 的顺序增长。请将其与二分法（每一步将区间对半分）比较，
        并用一句话说明为何
        计算器、GPS 接收机与神经网络训练
        都基于牛顿的思想。
\end{enumerate}

\medskip
\noindent\textbf{第四部分 --- 最坚固的梁。}
一根宽为 $w$、高度为 $d$ 的矩形梁，锯自一根
直径为 $30$ 厘米的圆木，故 $w^2 + d^2 = 900$。梁的
刚度与 $w d^2$ 成正比（高度要计两次：
搁栅是竖着放的！）。
\begin{enumerate}[resume]
  \item 将刚度表示为仅关于 $w$ 的函数：
        $S(w) = 900w - w^3$，定义在哪个区间上？
  \item 求导，列出变化表，并求出
        最优的 $w$ 与 $d$（先给出精确值，再精确到
        毫米）。
  \item 证明最优比例满足
        $d = w\sqrt2$ --- 古老木匠的比例。（附加趣闻：
        将圆木的直径三等分，
        并过分点作圆的垂线，
        恰好构造出这一矩形。）
  \item 最朴素的选择是\emph{正方形}梁
        （$w = d$）。计算它的刚度以及最优梁的刚度，
        并以百分比给出木匠的增益。
  \item 结语 --- 切线的三种作用，每种用一句话说明：
        几何（镜面定理）、数值计算（牛顿的
        “机器”）、最优化（梁），三者都借助同一个
        对象。展望下一步：十二年级将引入
        凸性 --- 即曲线位于切线的哪一侧 --- 用以
        深化上述三种应用。
\end{enumerate}
\end{problem}
